\chapter*{前言}
\markboth{前言}{前言}
\addcontentsline{toc}{chapter}{前言}
\sourcepage{vii}
本书源于作者在慕尼黑大学、雷根斯堡大学和明斯特大学讲授的黎曼曲面课程。本书旨在介绍这一丰富而优美的学科，同时介绍复流形理论中的一些方法；这些方法在单复变这一特殊情形下，显得格外初等而清晰。

全书分为三章。第一章把黎曼曲面作为覆盖空间来研究，并为此介绍所需的少量拓扑学基础。随后，我们构造由函数芽的解析延拓产生的黎曼曲面，其中尤其包括代数函数的黎曼曲面。我们还较详细地研究具有特殊多值性质的解析函数，例如全纯1形式的原函数和线性微分方程的解。

第二章讨论紧黎曼曲面，介绍Riemann--Roch定理、Abel定理及Jacobi反演问题等经典主要结果。层上同调是一项重要的技术工具。不过，我们只使用第一上同调群，而这类群处理起来相对容易。依照Serre的方法，所有主要定理都从以全纯函数层为系数的第一上同调群的有限维性推导出来。这一有限维性的证明，则以非齐次Cauchy--Riemann方程的局部可解性和Schwarz引理为基础。

第三章证明单连通黎曼曲面的Riemann映射定理（亦即单值化定理），以及Behnke--Stein关于非紧黎曼曲面的主要定理，即Runge逼近定理、Mittag--Leffler定理与Weierstrass定理。为此，我们使用Perron对Dirichlet问题的解法，
\sourcepage{viii}
以及Malgrange以Weyl引理为基础证明Runge逼近定理的方法。在这一章，我们还完成了第一章开始讨论的Stein定理，即具有指定自守加项的全纯函数的存在性定理，并介绍Röhrl对非紧黎曼曲面上Riemann--Hilbert问题的解法。

我们力求把预备知识的要求降到最低，并在需要时逐步建立必要的工具。不过，仍假定读者熟悉通常各用一学期讲授的单复变函数、一般拓扑学及代数学课程的内容。除此之外，第二、三章还用到微分拓扑和泛函分析中的少量事实，这些事实集中列在附录中。本书不需要Lebesgue积分，因为只对全纯函数或可微函数（相应地，对微分形式）作积分。我们也避免不加证明地引用任何关于曲面拓扑的定理。

本书内容大致相当于三个学期的课程。不过，第二、三章只以前面各章中的部分内容为基础。因此，读者读过第1、6、9节（黎曼曲面、层和微分形式的定义）后，可以直接阅读第二章。而在第二章中，只需第12至14节的内容，就足以在第三章中处理非紧黎曼曲面的主要定理。

英文版在每节末尾增添了习题，并在第8、17、29节中补充了一些段落。此外，有关覆盖的术语也作了调整。许多认真阅读德文版的读者帮助我们消除了若干错误，谨向他们致谢；特别感谢G.~Elencwajg，他还提出了部分习题。最后，同样要感谢译者B.~Gilligan的辛勤付出。

\begin{flushright}O. FORSTER\\明斯特，1981年5月\end{flushright}
\section*{第四次修订印刷补记}
在第二次和第四次印刷时，我们改正了若干排印错误及其他错误。谨感谢B.~Gilligan、B.~Elsner和O.~Hien编制勘误表。
\begin{flushright}O. FORSTER\\1999年4月\end{flushright}
